\chapter{Related Work}

\section{LLM-Based Role-Playing and Persona Simulation}

The use of LLMs as role-playing agents has progressed rapidly from single-turn character impersonation to sustained multi-turn persona maintenance. Park~\etal~\cite{park2023generative} introduced Generative Agents, demonstrating that 25 LLM-powered agents inhabiting a virtual town could exhibit emergent social behaviors such as party planning and relationship formation. This foundational work established the feasibility of LLM-based social simulation but focused on behavioral plausibility rather than quantitative fidelity to assigned personality profiles.

Subsequent work has addressed personality consistency more directly. Abdulhai~\etal~\cite{abdulhai2025} proposed using multi-turn reinforcement learning to maintain consistent persona behavior across extended conversations, demonstrating that RL-based fine-tuning can reduce persona drift relative to prompting-based approaches. Liu~\etal~\cite{liu2022extending} explored persona extending techniques that enrich a character's personality description with inferred attributes to improve consistency in dialogue systems.

The digital twin paradigm has emerged as a parallel line of research. Du~\etal~\cite{du2025twinvoice} proposed TwinVoice, a multi-dimensional benchmark for evaluating how well LLM personas replicate human behavior across personality, values, and decision-making dimensions. Li~\etal~\cite{li2025digitaltwin} examined how far LLMs are from serving as digital twins by benchmarking persona-based behavior chain simulation. Paglieri~\etal~\cite{paglieri2026persona} from Google DeepMind introduced Persona Generators, a framework for creating diverse synthetic personas at scale.

However, these works predominantly address single-agent persona fidelity---the question of whether one LLM agent can consistently impersonate one person. The qualitatively different challenge of maintaining persona fidelity when \textit{multiple agents interact}, exert social pressure on each other, and must sustain distinctive viewpoints against sycophantic convergence forces remains largely unexplored.


\section{Multi-Agent Systems for Deliberation}

Multi-agent LLM systems have been deployed for a range of collaborative and adversarial tasks. Frameworks such as AutoGen (Microsoft), CrewAI, and LangGraph provide infrastructure for orchestrating multi-agent workflows, but these are general-purpose tools without specific mechanisms for maintaining behavioral fidelity in deliberative settings.

In the policy domain, AgentSociety (Tsinghua University, 2025) demonstrated a simulation of 10,000 LLM-powered agents interacting over 5 million exchanges to model policy effects at population scale. While impressive in scale, this work does not evaluate individual agent behavioral fidelity---whether each agent's personality and argumentative style remain consistent with their assigned profile across interactions.

The Habermas Machine~\cite{bakker2024habermas}, published in \textit{Science}, represents a significant advance in AI-mediated deliberation, demonstrating that an AI system can generate consensus statements preferred by human participants over those produced by human mediators ($N=5{,}734$, 56\% preference). However, this system facilitates deliberation among \textit{real human} participants rather than simulating synthetic stakeholders, placing it in a fundamentally different problem space.

PolicySimEval (February 2025) represents the first benchmark specifically designed for evaluating agent-based policy simulations, comprising 20 comprehensive scenarios with 65 targeted sub-tasks and 200 auto-generated tasks. The best-performing systems achieve only 24.5\% accuracy, demonstrating the substantial difficulty of the policy simulation task and the massive room for improvement in current approaches.

The MAPS framework (Multi-agent AI-augmented Participatory System, 2025) uses DeepSeek~V3 for multi-round stakeholder deliberation, representing the closest existing work to the present study. However, MAPS does not include systematic behavioral fidelity evaluation---it focuses on participatory process outcomes rather than individual agent persona consistency.

Laban~\etal~\cite{laban2025lost} demonstrated that LLMs systematically degrade in performance during multi-turn conversations, a finding directly relevant to the persona drift problem addressed in this study. Shapira~\etal~\cite{shapira2026chaos} examined emergent dynamics in multi-agent LLM interactions, while Hu~\etal~\cite{hu2025dialoglab} presented DialogLab for authoring and simulating dynamic human-AI group conversations.

The gap in the literature is clear: while multi-agent deliberation systems exist at both population scale (AgentSociety) and human-mediated scale (Habermas Machine), and a benchmark has been established (PolicySimEval), no existing work addresses the specific problem of maintaining individual behavioral fidelity in a small-group multi-stakeholder deliberation among synthetic agents.


\section{OCEAN Personality Model in NLP}

The Big Five (OCEAN) personality model---comprising Openness, Conscientiousness, Extraversion, Agreeableness, and Neuroticism---has been the dominant framework in personality psychology since Costa and McCrae's (1992) formalization~\cite{costa1992neo}, with the Big Five Inventory (BFI) serving as the standard measurement instrument~\cite{john1999bigfive}.

The application of OCEAN to LLM behavior has generated substantial research interest. Lee~\etal~\cite{lee2025trait} developed TRAIT, a personality test set designed specifically for LLMs with psychometric rigor, finding that LLMs do exhibit distinct and partially consistent personality profiles. Hilliard~\etal~\cite{hilliard2024eliciting} investigated methods for eliciting personality traits in LLMs, examining how prompting strategies affect the manifestation of Big Five characteristics. Wang~\etal~\cite{wang2025psi} proposed a Personality Structured Interview framework for LLM simulation in personality research.

Recent work has begun connecting Big Five traits to LLM agent behavior in interactive settings. Research using the Sotopia testbed found that Agreeableness and Extraversion significantly affect believability, goal achievement, and knowledge acquisition in negotiations---a finding consistent with the present study's per-trait results. ``The Impact of Big Five Personality Traits on AI Agent Decision-Making in Public Spaces'' (March 2025) demonstrated that LLMs can simulate Big Five traits that influence collective decision-making, while ``Dynamic Personality in LLM Agents'' (ACL Findings 2025) proposed a framework for personality trait evolution in multi-agent simulations.

Despite these advances, the application of OCEAN to multi-stakeholder policy deliberation---where personality traits map to \textit{communication styles} rather than self-reported attitudes---remains unexplored. The present study's contribution is to operationalize OCEAN as a behavioral fidelity framework for deliberation simulation, using 30 linguistic features to continuously estimate trait expression and detect drift from assigned profiles.


\section{Behavioral Fidelity Evaluation}

The evaluation of behavioral fidelity in LLM-based simulation has emerged as a critical but underdeveloped research area. Gupta and Sheikh~\cite{gupta2026digital} demonstrated that LLM-powered social digital twins can simulate population behavioral responses to policy interventions with 20.7\% improvement over gradient boosting baselines, but their evaluation operates at the population level rather than the individual persona level.

The challenge of evaluation methodology itself has received attention. Stureborg~\etal~\cite{stureborg2024inconsistent} demonstrated that LLMs are inconsistent and biased when used as evaluators, while Verga~\etal~\cite{verga2024jury} proposed replacing single-model judges with jury panels of diverse models to improve evaluation reliability. These findings informed the present study's decision to use deterministic, rule-based behavioral feature extraction rather than LLM-based evaluation for measuring persona fidelity.

PolicySimEval (February 2025) provides the first benchmark for evaluating agent-based policy simulations, with best-performing systems achieving only 24.5\% accuracy, highlighting both the difficulty of the task and the nascent state of the field.

The gap in existing evaluation approaches is their focus on aggregate prediction accuracy---whether a simulated population votes correctly or a simulated consumer makes the expected purchase. No existing work measures \textit{turn-by-turn behavioral fidelity at the individual agent level} during multi-party interaction. This study addresses this gap through a 30-feature behavioral extraction pipeline that estimates OCEAN trait expression at each dialogue turn and quantifies drift from assigned personality priors.


\section{Existing Policy Simulation Tools}

The landscape of policy simulation tools can be categorized into four tiers by their proximity to the multi-stakeholder deliberation problem addressed in this study.

\textbf{Quantitative policy models} such as EUROMOD (EU Joint Research Centre) and the tools maintained by the U.S.\ Congressional Budget Office provide microsimulation of tax-benefit systems and macroeconomic projections. These tools are indispensable for fiscal analysis but do not model stakeholder behavior or communication dynamics.

\textbf{Agent-based modeling (ABM) platforms} such as AnyLogic (\$12,000--\$19,000/year license) and NetLogo (open-source) enable simulation of large agent populations with rule-based behaviors. Traditional ABM agents follow programmed rules and lack the linguistic richness and emergent reasoning capabilities of LLM-based agents. AgentTorch (MIT) bridges this gap partially by integrating LLMs into an ABM framework for large-scale social simulation.

\textbf{Commercial AI simulation platforms} represent the nearest competitive landscape. Simile (Stanford spin-off, \$100M Series~A) offers LLM-based persona replication for applications including earnings call rehearsal and litigation preparation. Aaru ($\sim$\$1B valuation) specializes in synthetic population prediction, achieving notable accuracy in election forecasting. Artificial Societies (YC W25, \EUR{}4.5M) provides 300--5,000 AI personas for marketing and communication testing at \$40/month. However, none of these platforms specifically targets small-group policy deliberation with behavioral fidelity guarantees. The nearest government-sector product is Fujitsu's Policy Twin (December 2024), which creates digital twins of societal movements for local government policy simulation in Japan; field trials identified policies that doubled both cost savings and health improvements.

\textbf{Citizen participation platforms} such as Polis (deployed in Taiwan's vTaiwan initiative, where 80\% of issues led to government action) and Stanford HAI's Deliberation.io (deployed by Washington, D.C.\ in September 2025) facilitate real human deliberation rather than synthetic simulation. DemocracyNext and MIT's Center for Constructive Communication launched a two-year lab in 2025 to integrate AI into citizens' assemblies. The UK's AI Consult tool automated analysis of public consultation responses, achieving a median review time of 23~seconds per response and saving over \pounds{}250,000 across seven live consultations.

The market gap is thus precisely defined: no existing tool simulates how 5--10 specific stakeholders would negotiate around a policy table, maintaining their distinctive argumentative styles while responding authentically to each other's positions.
